\documentclass[journal=jacsat,manuscript=article]{achemso}
\usepackage{graphicx}
\usepackage{booktabs}
\usepackage[version=4]{mhchem}

\author{First Author}
\email{first@example.org}
\affiliation{Department of Chemistry, University, City, Country}
\author{Second Author}
\affiliation{Institute, City, Country}

\title{Title of the Paper}

\keywords{first keyword, second keyword}

\begin{document}

\begin{abstract}
State the problem, what you did, and the main result in a few sentences.
\end{abstract}

\section{Introduction}
\label{sec:intro}
Introduce the problem and why it matters, then the contributions of this
paper. Cite related work like this~\cite{kohn1965}, or several works at
once~\cite{kohn1965,becke1993}.

\section{Experimental Section}
\label{sec:method}
Describe what you did in enough detail to reproduce it. Numbered equations use
the \texttt{equation} environment:
\begin{equation}
  \label{eq:main}
  \ce{CO2 + 4 H2 -> CH4 + 2 H2O}
\end{equation}
Refer to it as Equation~\eqref{eq:main}.

\section{Results}
\label{sec:results}
Table~\ref{tab:results} and Figure~\ref{fig:overview} show the main results.

\begin{table}[htbp]
  \centering
  \caption{Computed bond lengths.}
  \label{tab:results}
  \begin{tabular}{lrr}
    \toprule
    Functional & Bond length (\AA) & Error (\AA) \\
    \midrule
    B3LYP & 1.21 & 0.03 \\
    PBE0 & \textbf{1.18} & 0.02 \\
    \bottomrule
  \end{tabular}
\end{table}

\begin{figure}[htbp]
  \centering
  \includegraphics[width=0.6\linewidth]{example-image}
  \caption{Overview of the method. Upload an image and use its file name here.}
  \label{fig:overview}
\end{figure}

\section{Discussion}
Interpret the results, compare them with earlier work and state the limitations.

\section{Conclusions}
Summarise the findings and what comes next.

\begin{acknowledgement}
Thank the people and funding that supported this work.
\end{acknowledgement}

\begin{suppinfo}
List the Supporting Information: methods, spectra, coordinates.
\end{suppinfo}

\bibliography{references}

\end{document}
